\chapter{Linar Maps}

\lecture{7.1}{}{Linear Maps}
\section{Linear Maps}

\begin{definition}[Linear Transformation]\label{def:linear-transformation} Let
$F$ be a field. A function \[ T: F^n \to F^m \] is called a \textbf{linear
transformation} if: \begin{enumerate}[label=(\roman*)] \item $\forall u, w \in
F^n$, $T(u + w) = T(u) + T(w)$ \item $\forall u \in F^n$, $\forall a \in F$,
$T(au) = aT(u)$ \end{enumerate} \end{definition}

\begin{note} It is not required that $n = m$, but the field $F$ must be the
same for both vector spaces. \end{note}

\begin{theorem}\label{thm:linear-map-zero} If $T: F^n \to F^m$ is a linear map,
then: \begin{enumerate}[label=(\roman*)] \item $T(\underline{0}) =
\underline{0}$ \item $T(-u) = -T(u)$ \end{enumerate} \end{theorem}

\begin{proof} Recall that $\forall u \in F^n$, $0u = \underline{0}$ and $(-1)u
= -u$. \begin{enumerate}[label=(\roman*)] \item $T(\underline{0}) =
T(0\underline{0}) = 0T(\underline{0}) = \underline{0}$ \item $T(-u) = T((-1)u) =
(-1)T(u) = -T(u)$ \end{enumerate} \end{proof}

\begin{note} Theorem~\ref{thm:linear-map-zero} describes necessary but not
sufficient conditions for a mapping $T$ to be linear. Use this theorem to
quickly rule out linearity. \end{note}

\begin{theorem}[Linear Combination Preservation]\label{thm:
linear-combination-preservation} Let $T: F^n \to F^m$ be a linear transformation.
For any $u_1, \ldots, u_k \in F^n$ and any $c_1, \ldots, c_k \in F$, \[
T\left(\sum_{i=1}^{k} c_i u_i\right) = \sum_{i=1}^{k} c_i T(u_i) \]
\end{theorem}

\begin{proof} By induction on $k$. Base case ($k = 1$): $T(c_1 u_1) = c_1
T(u_1)$ is the definition of a linear transformation. Induction hypothesis:
Assume $T\left(\sum_{i=1}^{k-1} c_i u_i\right) = \sum_{i=1}^{k-1} c_i T(u_i)$.
Induction step: \begin{align*} T\left(\sum_{i=1}^{k} c_i u_i\right) &=
T\left(\sum_{i=1}^{k-1} c_i u_i + c_k u_k\right) = T\left(\sum_{i=1}^{k-1}
c_i u_i\right) + c_k T(u_k) \\ &= \sum_{i=1}^{k-1} c_i T(u_i) + c_k T(u_k) =
\sum_{i=1}^{k} c_i T(u_i) \end{align*} \end{proof}

\lecture{7.2}{}{Map Arithmetic} \section{Map Arithmetic}

\begin{definition}[Sum of Linear Maps]\label{def:sum-linear-maps} Let $T, S$ be
linear maps from $F^n$ to $F^m$. The sum $(S + T)$ is the function defined by \[
(S + T)(u) \coloneqq S(u) + T(u) \] \end{definition}

\begin{theorem}\label{thm:sum-is-linear} Let $T, S$ be linear maps from $F^n$
to $F^m$. Then $T + S: F^n \to F^m$ is also a linear map. \end{theorem}

\begin{definition}[Scalar Multiplication of Linear Maps]\label{def:
scalar-mult-linear-maps} Let $T: F^n \to F^m$ be a linear map and let $c \in F$
be a scalar. The multiplication of $T$ by the scalar $c$, denoted $cT$, is the
function defined by \[ (cT)(u) \coloneqq c(T(u)) \] \end{definition}

\begin{theorem}\label{thm:scalar-mult-is-linear} Let $T: F^n \to F^m$ be a
linear map. Then $cT: F^n \to F^m$ is also a linear map. \end{theorem}

\begin{theorem}[Properties of Linear Map Addition]\label{thm:
linear-map-addition-properties} Let $T, S, W: F^n \to F^m$ be linear maps. Then:
\begin{enumerate}[label=(\roman*)] \item $T + S = S + T$ \item $T + (S + W) = (T
+ S) + W$ \item The map $O: F^n \to F^m$ defined by $O(u) = \underline{0}$
satisfies $T + O = T$ \item The map $(-1)T$, denoted $-T$, satisfies $T + (-T) =
O$ \end{enumerate} \end{theorem}

\begin{definition*}[Composition/Product of Linear Maps]\label{def:
composition-linear-maps} Let $T: F^n \to F^m$ and $S: F^m \to F^p$ be linear
maps. The composition $S \circ T: F^n \to F^p$ is defined by \[ (S \circ T)(u) =
S(T(u)) \] In the context of linear algebra, this operation is also called the
\textbf{product} of the maps and denoted $ST \coloneqq S \circ T$.
\end{definition*}

\begin{theorem}\label{thm:composition-is-linear} Let $T: F^n \to F^m$ and $S:
F^m \to F^p$ be linear maps. Then the composition $S \circ T: F^n \to F^p$ is
also a linear map. \end{theorem}

\begin{note} The product of linear maps is associative (since function
composition is associative) but not commutative. \end{note}

\begin{definition*}[Inverse Map]\label{def:inverse-map} A linear map $T: F^n \to
F^n$ that is one-to-one and onto is called \textbf{invertible}. Its inverse
$T^{-1}$ satisfies \[ T^{-1}T = TT^{-1} = I_n \] where $I_n: F^n \to F^n$ is the
identity map defined by $I_n(u) = u$. \end{definition*}

\begin{theorem}\label{thm:inverse-is-linear} Let $T: F^n \to F^n$ be a linear
map. If $T$ is invertible, then its inverse map is linear. \end{theorem}

\begin{proof} $T$ is invertible if and only if it is one-to-one and onto. Let
$w_1, w_2 \in F^n$ and $c \in F$. Since $T$ is onto, $\exists u_1, u_2 \in F^n$
such that $T(u_1) = w_1$ and $T(u_2) = w_2$. Therefore $T^{-1}(w_1) = u_1$ and
$T^{-1}(w_2) = u_2$. Since $T$ is linear, $T(u_1 + u_2) = w_1 + w_2$. Therefore
$T^{-1}(w_1 + w_2) = u_1 + u_2 = T^{-1}(w_1) + T^{-1}(w_2)$. Since $T$ is linear,
$T(cu_1) = cT(u_1)$. Therefore $T^{-1}(cw_1) = cu_1 = cT^{-1}(w_1)$.
\end{proof}

\lecture{7.3}{}{Map Basisification} \section{Map Basisification}

\begin{theorem}[Image of Basis Determines the Map]\label{thm:
basis-determines-map} Let $T: F^n \to F^m$ and $S: F^n \to F^m$ be linear maps
and let $b_1, \ldots, b_n \in F^n$ be a basis of $F^n$. If $ T(b_i) = S(b_i)
\quad \forall i \in \{1, \ldots, n\} $ then \[ T(u) = S(u) \quad \forall u \in
F^n \] \end{theorem}

\begin{proof} Let $u \in F^n$. Since $b_1, \ldots, b_n$ are a basis, they span
$F^n$, so $u$ is a linear combination of them: \[ u = \alpha_1 b_1 + \cdots +
\alpha_n b_n \] Since $T, S$ are linear maps, \begin{align*} T(u) &= T(\alpha_1
b_1 + \cdots + \alpha_n b_n) \\ &= \alpha_1 T(b_1) + \cdots + \alpha_n T(b_n) \\
&= \alpha_1 S(b_1) + \cdots + \alpha_n S(b_n) \\ &= S(\alpha_1 b_1 + \cdots +
\alpha_n b_n) = S(u) \end{align*} \end{proof}

\begin{theorem}[Any Image of a Basis Makes a Linear Map]\label{thm:
basis-makes-map} Let $b_1, \ldots, b_n$ be a basis of $F^n$ and let $w_1, \ldots,
w_n$ be any $n$ vectors in $F^m$. There exists a unique linear map $T: F^n \to
F^m$ such that \[ T(b_i) = w_i, \quad \forall i \in \{1, \ldots, n\} \]
\end{theorem}

\begin{proof} Since $b_1, \ldots, b_n$ are a basis, any vector in $F^n$ is a
unique linear combination of the $b_i$. Define a transformation $T: F^n \to F^m$
that maps every $u = \alpha_1 b_1 + \cdots + \alpha_n b_n$ to \[ T(u) = \alpha_1
w_1 + \cdots + \alpha_n w_n \] If $T$ is linear, Theorem~\ref{thm:
basis-determines-map} proves it's unique. We prove that $T$ is linear. Let $u_1
= \alpha_1 b_1 + \cdots + \alpha_n b_n$ and $u_2 = \beta_1 b_1 + \cdots +
\beta_n b_n$, and $c \in F$. \begin{align*} T(u_1 + u_2) &= T((\alpha_1 b_1 +
\cdots + \alpha_n b_n) + (\beta_1 b_1 + \cdots + \beta_n b_n)) \\ &= T((\alpha_1
+ \beta_1)b_1 + \cdots + (\alpha_n + \beta_n)b_n) \\ &= (\alpha_1 + \beta_1)w_1
+ \cdots + (\alpha_n + \beta_n)w_n \\ &= (\alpha_1 w_1 + \cdots + \alpha_n w_n)
+ (\beta_1 w_1 + \cdots + \beta_n w_n) \\ &= T(u_1) + T(u_2) \end{align*}
Similarly, $T(cu_1) = cT(u_1)$. \end{proof}

\lecture{7.4}{}{Map Matrixification} \section{Map Matrixification}

\begin{note} A matrix of order $n \times 1$ is also called a column vector. We
use the notation $Ax = b$ to represent a system of linear equations, where $x$
and $b$ are column vectors. We also use the underbar to denote $n$-tuples in
$F^n$, and we also call $n$-tuples in $F^n$ vectors. There is no reason to
constantly be on guard about the difference between column vectors and
$n$-tuples. \end{note}

\begin{theorem}[Matrices Define Linear Maps]\label{thm:matrix-defines-map} Let
$A \in M_{m \times n}(F)$. The map $T_A: F^n \to F^m$ defined by \[ T(u) = Au \]
is a linear map. \end{theorem}

\begin{proof} Let $u_1, u_2 \in F^n$ and $c \in F$. $T(u_1 + u_2) = A(u_1 +
u_2) = Au_1 + Au_2 = T(u_1) + T(u_2)$ by properties of matrix multiplication.
$T(cu_1) = A(cu_1) = c(Au_1) = cT(u_1)$ by properties of matrix multiplication.
\end{proof}

\begin{theorem}[Standard Basis Matrix]\label{thm:standard-basis-matrix} Let
$e_1, \ldots, e_n$ be the standard basis of $F^n$ and let $w_1, \ldots, w_n \in
F^m$ be any $n$ vectors. Let $T: F^n \to F^m$ be the linear map defined by \[
T(e_i) = w_i \] and let $A \in M_{m \times n}(F)$ be the matrix whose entries
are \[ A_c^i = w_i \] (i.e., the matrix whose columns are $w_1, \ldots, w_n$).
Then, $\forall u \in F^n$, \[ T(u) = Au \] \end{theorem}

\begin{proof} Let $u = (\alpha_1, \ldots, \alpha_n) = \alpha_1 e_1 + \cdots +
\alpha_n e_n$ and denote by $w_{ij}$ the $i$-th component of $w_j$. By
Theorem~\ref{thm:linear-combination-preservation}, $T(u) = \alpha_1 w_1 + \cdots
+ \alpha_n w_n$. By properties of matrix algebra, \begin{align*} Au &=
A(\alpha_1 e_1 + \cdots + \alpha_n e_n) \\ &= \alpha_1 Ae_1 + \cdots + \alpha_n
Ae_n \\ &= \alpha_1 \begin{bmatrix} w_{11} \\ \vdots \\ w_{m1} \end{bmatrix} +
\cdots + \alpha_n \begin{bmatrix} w_{1n} \\ \vdots \\ w_{mn} \end{bmatrix} \\ &=
\alpha_1 w_1 + \cdots + \alpha_n w_n \end{align*} \end{proof}

\begin{notation} The standard-basis matrix of a map $T$ is denoted by $[T]$.
\end{notation}

\lecture{7.5}{}{Map Arithmetic the Easy Algorithmic Way} \section{Map
Arithmetic the Easy Algorithmic Way}

\begin{theorem}[Sum and Scalar Multiplication]\label{thm:sum-scalar-mbas} Let
$T: F^n \to F^m$ and $S: F^n \to F^m$ be linear maps, and $c \in F$. Then:
\begin{enumerate}[label=(\roman*)] \item $[T + S] = [T] + [S]$ \item $[cT] =
c[T]$ \end{enumerate} \end{theorem}

\begin{proof} By properties of matrix addition and scalar multiplication:
\begin{enumerate}[label=(\roman*)] \item $\forall u \in F^n$, $[T + S]u = [T]u +
[S]u = T(u) + S(u) = (T + S)(u)$ \item $\forall u \in F^n$, $c[T]u = c([T]u) =
c(T(u)) = (cT)(u)$ \end{enumerate} \end{proof}

\begin{theorem}[Product of Maps]\label{thm:product-mbas} Let $T: F^n \to F^m$
and $S: F^m \to F^p$ be linear maps. Then \[ [ST] = [S][T] \] \end{theorem}

\begin{proof} $\forall u \in F^n$, \begin{align*} [S][T]u &= [S]([T]u) \quad
\text{(by associativity of matrix multiplication)} \\ &= [S](T(u)) \quad
\text{(because $[T]$ is the standard-basis matrix of $T$)} \\ &= S(T(u)) \quad
\text{(because $[S]$ is the standard-basis matrix of $S$)} \\ &= (ST)(u) \quad
\text{(definition of $ST$)} \end{align*} \end{proof}

\begin{theorem}[Inverse]\label{thm:inverse-mbas} Let $T: F^n \to F^n$ be a
linear map. If $T$ is one-to-one and onto, then \[ [T^{-1}] = [T]^{-1} \]
\end{theorem}

\begin{proof} $\forall u \in F^n$, \begin{align*} [T^{-1}][T]u &=
[T^{-1}]([T]u) \\ &= [T^{-1}](T(u)) \\ &= T^{-1}(T(u)) \\ &= I_n u \\ &= u
\end{align*} \end{proof}

\begin{note} Finding a function $g$ such that $g(f(x)) = x$ for all $x$ does
not imply that $f$ is one-to-one and/or onto, not even when $f$ is a linear map.
\end{note}

\lecture{7.6}{}{Eigenvectors} \section{Eigenvectors}

\begin{definition}[Eigenvector and Eigenvalue]\label{def:eigenvector} Let $T:
F^n \to F^n$ be a linear map. A vector $u \neq \underline{0} \in F^n$ is called
an \textbf{eigenvector} of $T$ if there exists a scalar $\lambda \in F$ such
that \[ T(u) = \lambda u \] The scalar $\lambda$ is then called an
\textbf{eigenvalue}. \end{definition}

\begin{note} The vector $\underline{0}$ doesn't count as an eigenvector (it is
always true that $T(\underline{0}) = \underline{0} = \lambda \underline{0}$, for
any $T$, for any $\lambda$). \end{note}

\begin{theorem}\label{thm:eigenvector-scalar-multiple} Let $T: F^n \to F^n$ be
a linear map. If $u$ is an eigenvector of $T$ belonging to the eigenvalue
$\lambda$, then so is $w = cu$, for any $c \neq 0 \in F$. \end{theorem}

\begin{proof} \[ T(w) = T(cu) = cT(u) = c\lambda u = \lambda cu = \lambda w \]
\end{proof}

\lecture{7.7}{}{How to Find Eigenvalues and Eigenvectors} \section{How to Find
Eigenvalues and Eigenvectors}

\begin{note} Using map matrixification (non-canonical), we can describe an
algorithm for finding eigenvalues and eigenvectors that sometimes works.
\end{note}

\begin{theorem*}[Matrixification of the Eigenproblem]\label{thm:
eigenproblem-matrix} Let $T: F^n \to F^n$ be a linear map and let $[T]$ denote
the standard-basis matrix of the map $T$. We seek a vector $u \neq \underline{0}
\in F^n$ such that $T(u) = \lambda u$ for some $\lambda \in F$. This is
equivalent to seeking a vector $u \neq \underline{0} \in M_{n \times 1}(F)$ such
that $[T]u = \lambda u$. The following are equivalent:
\begin{enumerate}[label=(\roman*)] \item $T(u) = \lambda u$ \item $[T]u =
\lambda u$ \item $\lambda u - [T]u = \underline{0}$ \item $\lambda I_n u - [T]u
= \underline{0}$ \item $(\lambda I_n - [T])u = \underline{0}$ \end{enumerate}
Thus, we seek a non-trivial solution to the homogeneous system of equations
whose (not-augmented) coefficient matrix is $\lambda I_n - [T]$, for some
$\lambda \in F$. \end{theorem*}

\begin{note} A non-trivial solution to the system $(\lambda I_n - [T])u =
\underline{0}$ exists if and only if $\det(\lambda I_n - [T]) = 0$. \end{note}

\begin{theorem}[Characteristic Polynomial]\label{thm:characteristic-polynomial}
Let $T \in M_n(F)$ and $x$ a variable. The expression \[ \det(xI_n - T) \] is a
polynomial of degree $n$ in the variable $x$. \end{theorem}

\begin{proof} By induction on $n$. \end{proof}

\begin{definition*}[Characteristic Polynomial and Equation]\label{def:
characteristic} The polynomial $\det(xI_n - [T])$ is called the
\textbf{characteristic polynomial} of $T$. The equation $\det(xI_n - [T]) = 0$
is called the \textbf{characteristic equation} of $T$. \end{definition*}

\begin{conclusion*}\label{conc:eigenvalues-roots} Let $T: F^n \to F^n$ be a
linear map. The eigenvalues of $T$ are the roots of the characteristic
polynomial of $T$. They can be found by solving the characteristic equation of
$T$. \end{conclusion*}

\begin{note} We transformed the problem of finding eigenvalues into the problem
of finding roots of polynomials, which is not an easy task. For quadratic
polynomials, standard formulas exist. For cubic polynomials, formulas exist but
can be difficult to use. For quartic and higher-degree polynomials, no general
algebraic solution exists. Numerical algorithms exist for finding eigenvalues
efficiently, but they do not rely on solving the characteristic equation.
\end{note}

\begin{note} The characteristic equation remains useful for theoretical
purposes and for hand calculations on small matrices. \end{note}